\section{模型假设及适用范围}

题设参数按确定值处理。无人机沿两节点水平直线飞行，巡航高度由DEM净空决定；不考虑风场、临时禁飞区与地面道路通行。不可拆货箱是运输的最小单位，一箱只能交付一次。第一问每架次仅服务一个服务区；后两问允许同架次多点交付，并按交付后的剩余箱重更新能耗。

准备和装载从架次开始时占用实体机及电池；返航结束后实体机可开始下一次准备，电池须先充满。附件未规定共享工位数量，因此不额外施加工位互斥。中继在给定悬停点保持不动。每次任务需由O01出发并返回，服务前完成准备和建链；悬停和通信附加功率依附件计入能源消耗。中继任务结束与下一次准备之间保留给定的300~s周转时间。

以完整交接完成时刻检验医疗物资期望时刻及首批保障截止时刻。其他货箱的期望时刻只形成软性迟到惩罚。这个优先级与题面“应满足”和“用于衡量”的措辞一致。题设DEM按地形表面使用；其标称高程、节点高程与通信设备高度视为同一垂直基准。若存在较大定位误差或未表示的新障碍，需重新核算最小通信裕量。

关于目标权衡，本文不把单位不同的指标直接相加。第一问采用字典序；第二、三问先满足所有硬约束，再以优先系数加权期望迟到为主指标，比较总完成时间、能耗和架次数。由于后两问通过候选生成及有限时求解得到，所谓“优化方案”指相对可行基线的改进，不代表全局最优。

